% GENERATED FILE. Edit canonical lessons, then run npm run generate:books.
% canonical-source-hash: fnv1a64:958e3d97c4652da6
% canonical-lessons: ES-C749-gemelo, ES-C749-hermanastro, ES-C749-adolescente, ES-C749-bisabuelo, ES-C749-huerfano

\chapter{Más parientes}
\label{ch:es-mas-parientes}
\hlchaptermodality{749}

\begin{chapteropening}
\textbf{By the end of this chapter:} \emph{I can describe more family members: a teenager, twins, a stepbrother, a great-grandfather and an orphan.}

Complete the last lesson of chapter 749: I can describe more family members: a teenager, twins, a stepbrother, a great-grandfather and an orphan.
\end{chapteropening}

\section[el gemelo]{el gemelo — the twin; the calf (muscle)}

\label{lesson:ES-C749-gemelo}



\begin{quote}
Before the new one: say the Spanish for the stepmother, then the Spanish for the stepfather.
\end{quote}

\subsection*{You'll want to know: el gemelo}
\textbf{el gemelo} — \textquotedblleft{}the twin; the calf (muscle)\textquotedblright{}.

\begin{quote}
Mi hermana tiene dos gemelos.
\end{quote}

\begin{quote}
\emph{My sister has twins.}
\end{quote}

\emph{Los gemelos} are also the calf muscles and cufflinks. Non-identical twins are \emph{mellizos}.

\begin{cousinweb}
\emph{gemelo} is a learned borrowing from Latin \emph{gemellus} \textquotedblleft{}twin\textquotedblright{} (EtymDB), a small form of \emph{geminus} \textquotedblleft{}twin\textquotedblright{}.
\end{cousinweb}

\begin{grammarlens}[title={What you've built}]
One more word for telling, explaining and describing at B1.
\end{grammarlens}

\subsection*{Your turn}
\begin{itemize}
\raggedright
  \item \emph{Say it:} \emph{el gemelo}
  \item \emph{Say it:} \emph{el gemelo}, once more
  \item \emph{Say it:} a sentence of your own with \emph{el gemelo}
  \item \emph{Recall:} say the Spanish for the stepmother, then the Spanish for the stepfather, then say \emph{el gemelo} again
\end{itemize}

\begin{tcolorbox}[breakable,colback=teal!4,colframe=teal!35!black,title={Before you move on}]
What does el gemelo mean? (\textquotedblleft{}The twin; the calf (muscle)\textquotedblright{}.) Say it once more.
\end{tcolorbox}
\section[el hermanastro]{el hermanastro — the stepbrother}

\label{lesson:ES-C749-hermanastro}



\begin{quote}
Before the new one: say the Spanish for the stepfather, then the Spanish for the twin.
\end{quote}

\subsection*{You'll want to know: el hermanastro}
\textbf{el hermanastro} — \textquotedblleft{}the stepbrother\textquotedblright{}.

\begin{quote}
Mi hermanastro vive con nosotros.
\end{quote}

\begin{quote}
\emph{My stepbrother lives with us.}
\end{quote}

\emph{La hermanastra} — stepsister.

\begin{cousinweb}
\emph{hermanastro} is built on \emph{hermano} with the ending \emph{-astro} (EtymDB); \emph{hermano} is inherited from Latin \emph{germānus} \textquotedblleft{}of the same parents\textquotedblright{}.
\end{cousinweb}

\begin{grammarlens}[title={What you've built}]
One more word for telling, explaining and describing at B1.
\end{grammarlens}

\subsection*{Your turn}
\begin{itemize}
\raggedright
  \item \emph{Say it:} \emph{el hermanastro}
  \item \emph{Say it:} \emph{el hermanastro}, once more
  \item \emph{Say it:} a sentence of your own with \emph{el hermanastro}
  \item \emph{Recall:} say the Spanish for the stepfather, then the Spanish for the twin, then say \emph{el hermanastro} again
\end{itemize}

\begin{tcolorbox}[breakable,colback=teal!4,colframe=teal!35!black,title={Before you move on}]
What does el hermanastro mean? (\textquotedblleft{}The stepbrother\textquotedblright{}.) Say it once more.
\end{tcolorbox}
\section[el adolescente]{el adolescente — the teenager}

\label{lesson:ES-C749-adolescente}



\begin{quote}
Before the new one: say the Spanish for the twin, then the Spanish for the stepbrother.
\end{quote}

\subsection*{You'll want to know: el adolescente}
\textbf{el adolescente} — \textquotedblleft{}the teenager\textquotedblright{}.

\begin{quote}
Mi hijo es un adolescente.
\end{quote}

\begin{quote}
\emph{My son is a teenager.}
\end{quote}

\emph{La adolescente} for a girl.

\begin{cousinweb}
\emph{adolescente} is a learned borrowing from Latin \emph{adolēscēns} \textquotedblleft{}growing up\textquotedblright{} (EtymDB), the participle of \emph{adolēscere}.
\end{cousinweb}

\begin{grammarlens}[title={What you've built}]
One more word for telling, explaining and describing at B1.
\end{grammarlens}

\subsection*{Your turn}
\begin{itemize}
\raggedright
  \item \emph{Say it:} \emph{el adolescente}
  \item \emph{Say it:} \emph{el adolescente}, once more
  \item \emph{Say it:} a sentence of your own with \emph{el adolescente}
  \item \emph{Recall:} say the Spanish for the twin, then the Spanish for the stepbrother, then say \emph{el adolescente} again
\end{itemize}

\begin{tcolorbox}[breakable,colback=teal!4,colframe=teal!35!black,title={Before you move on}]
What does el adolescente mean? (\textquotedblleft{}The teenager\textquotedblright{}.) Say it once more.
\end{tcolorbox}
\section[el bisabuelo]{el bisabuelo — the great-grandfather}

\label{lesson:ES-C749-bisabuelo}



\begin{quote}
Before the new one: say the Spanish for the stepbrother, then the Spanish for the teenager.
\end{quote}

\subsection*{You'll want to know: el bisabuelo}
\textbf{el bisabuelo} — \textquotedblleft{}the great-grandfather\textquotedblright{}.

\begin{quote}
Mi bisabuelo vivió cien años.
\end{quote}

\begin{quote}
\emph{My great-grandfather lived to a hundred.}
\end{quote}

\emph{Los bisabuelos} — great-grandparents.

\begin{cousinweb}
\emph{bisabuelo} is built on Latin \emph{bis} \textquotedblleft{}twice\textquotedblright{} and \emph{abuelo} \textquotedblleft{}grandfather\textquotedblright{}, which is usually traced to Latin \emph{aviolus}, from \emph{avus} \textquotedblleft{}grandfather\textquotedblright{}.
\end{cousinweb}

\begin{grammarlens}[title={What you've built}]
One more word for telling, explaining and describing at B1.
\end{grammarlens}

\subsection*{Your turn}
\begin{itemize}
\raggedright
  \item \emph{Say it:} \emph{el bisabuelo}
  \item \emph{Say it:} \emph{el bisabuelo}, once more
  \item \emph{Say it:} a sentence of your own with \emph{el bisabuelo}
  \item \emph{Recall:} say the Spanish for the stepbrother, then the Spanish for the teenager, then say \emph{el bisabuelo} again
\end{itemize}

\begin{tcolorbox}[breakable,colback=teal!4,colframe=teal!35!black,title={Before you move on}]
What does el bisabuelo mean? (\textquotedblleft{}The great-grandfather\textquotedblright{}.) Say it once more.
\end{tcolorbox}
\section[el huérfano]{el huérfano — the orphan}

\label{lesson:ES-C749-huerfano}



\begin{quote}
Before the new one: say the Spanish for the teenager, then the Spanish for the great-grandfather.
\end{quote}

\subsection*{You'll want to know: el huérfano}
\textbf{el huérfano} — \textquotedblleft{}the orphan\textquotedblright{}.

\begin{quote}
Se quedó huérfano a los diez años.
\end{quote}

\begin{quote}
\emph{He was orphaned at ten.}
\end{quote}

\emph{Quedarse huérfano} — to be orphaned.

\begin{cousinweb}
\emph{huérfano} is usually traced to Late Latin \emph{orphanus}, borrowed from Greek \emph{orphanós} \textquotedblleft{}bereft\textquotedblright{}.
\end{cousinweb}

\begin{grammarlens}[title={What you've built}]
One more word for telling, explaining and describing at B1.
\end{grammarlens}

\subsection*{Your turn}
\begin{itemize}
\raggedright
  \item \emph{Say it:} \emph{el huérfano}
  \item \emph{Say it:} \emph{el huérfano}, once more
  \item \emph{Say it:} a sentence of your own with \emph{el huérfano}
  \item \emph{Recall:} say the Spanish for the teenager, then the Spanish for the great-grandfather, then say \emph{el huérfano} again
\end{itemize}

\begin{tcolorbox}[breakable,colback=teal!4,colframe=teal!35!black,title={Before you move on}]
What does el huérfano mean? (\textquotedblleft{}The orphan\textquotedblright{}.) Say it once more.
\end{tcolorbox}
